\section{The September continuation}\label{sec:items}

The continuation of 17--23 September is organised as the 28-item \emph{Start reading} list of the README, with proof packages in \texttt{research/incoming/}. The workbench describes the results of its 22 September release as written proofs with Python replays, without a Lean build or independent human review (\note{43a\_} §2, citing \texttt{DAILY\_RESULTS\_20260922}).

\subsection{The audited items}

The 28 items of the README's \emph{Start reading} list (17--23 September) are listed in \note{43a\_} §2, with their packages. Those audited:
\begin{itemize}
\item \emph{Item 6} (audited). With $L_a=8\tau(a)^2/\varphi(R_a)-\mathcal E_a$, where $\mathcal E_a$ is the residue-collision energy of the $2\tau(a)$ divisors of $pa$ modulo $R_a$: for $p\equiv1\pmod4$, $p\ge2^{20}$, $\sum_{a\in(p/4,p/2)}L_a\le-\tfrac1{24}p\log p$. So the lower bound that keeps only the principal character, summed over the shells, is negative for large $p$, and that bound alone does not show occupancy.
\item \emph{Item 7} (audited). For $p\equiv1\pmod 8$, bounds in terms of the least quadratic nonresidue; at the hard primes a middle state $(a,u)$ on a shell $p/4<a<p/2$ has $a\le3(p+3)/11$, and an exterior state satisfies $22(p+1)\le23(p-2a+1)^2$. The workbench says this localises solutions and does not prove the region occupied.
\item \emph{Item 14} (audited). For $h'=4j-1>t$ the divisors $W\mid j^2$ with $W\equiv t\pmod{h'}$ are canonical, companion or finite templates; the finite ones satisfy the sharp bound $h'\le t^2-3t+1$, with equality exactly in one parametrised family, attained at infinitely many primes $p\equiv1\pmod{840}$. For every $t\ge4$ infinitely many hard primes have an exterior state whose same-grade middle box is empty.
\item \emph{Item 21} (audited). For a proper rational trace source the complete set of same-word residual returns is $\{k\mid R: d\mid k,\ k\equiv1\pmod{4K(u)}\}$; a uniform return exists exactly for the nine words $u\mid36$.
\item \emph{Item 22} (audited, certificate). For every solution at $p\equiv1\pmod{12}$ an explicit symmetric polynomial $\mathfrak N(p,x,y,z)$ of degree $8$ satisfies $\mathfrak N<-\frac{125873811}{262144}p^8$ ($\approx-480.17\,p^8$), and $\mathfrak N<-6960.19\,p^8$ when $p\equiv1\pmod{24}$; the proof reduces to finite polynomial positivity certificates, recomputed independently. It is a statement about existing solutions, not about their existence.
\item \emph{Item 26} (audited in part: the least example and the deduction from it; its prime-square trace theorem at $p\equiv1\pmod{24}$ and its terminality statement read and correct as read, with the identity $16(4q^3+1)-(4q+1)(16q^2-4q+1)=15$ checked; the factor-sum, Pell and fixed-tail-sum sections not checked). $p=67369$ is the least prime $p\equiv1\pmod{24}$ whose first \emph{trace shell} ($a=16849$, $R=27$) precedes its first full shell ($a=16850$, $R=31$); so no universal map from rational traces to solutions avoids increasing $a$.
\item \emph{Item 4} (conditional on a computer enumeration of $1{,}381{,}117{,}764$ rooted profiles, $5{,}922{,}259$ of them range-compatible, run by two implementations of the workbench and not re-run here). At every hard prime, from any nonresidue prime, factoring at most 25 numbers $(p+\sigma qt^2)/4$ with $t\in\{1,3,5,7,9\}$ yields at least six distinct nonresidue primes. A direct test of all 108 hard primes below $40{,}000$ agrees.
\end{itemize}
Items 2 and 3 are method counterexamples, and their certificates are audited: the seven-cycle sink of item~2 at $p=2521$ and the closed triangle of item~3; their closure and escape theorems are located. For item~8, the middle identity $p=RQ-(Q+1)^2/(4A)$ and its two bounds were reproduced on all $1{,}613$ oriented middle states at primes $p\equiv1\pmod8$ below $4000$, and its atlas is located. The $22$-digit prime $7510085481569082811681$ of item~3 is proved prime here by a Pocklington chain ($p-1=2^5\cdot3\cdot5\cdot q_1$, $q_1-1=2\cdot3\cdot5\cdot351707\cdot q_2$, $q_2-1=2\cdot3\cdot17^2\cdot5051\cdot169307$, witnesses below $12$ at every level, the leaves by trial division). Item 23 is audited in part (its certificates; the deduction read in part); items 11--13, 15 and 20 use imported material (Section~\ref{sec:overlaps}); the remaining 17 items, including these five, are located. The results bench R01--R10 and the arguments X1--X7 are checked in \note{43a\_} §3; one source passage error (X5: $289^3\equiv169$, not $-1$, modulo $840$) is recorded by the workbench itself.

\subsection{The \texorpdfstring{$(3,4,\infty)$}{(3,4,infinity)} monodromy covers}\label{ssec:es-s6}

A package of the workbench (ES-09) imports three integer matrices $A_1,A_2,A_\infty$ with $A_1^3=A_2^4=A_1A_2A_\infty=I$ from the $S^6$ manuscript \cite{AlpogeS6}, written out in the package's \texttt{src/core.tex}, and nothing else from it. They are the duals $A_j=(T_j^{-1})^{\mathsf T}$ of the local monodromies $T_1,T_2$ of the $S^6$ period system, acting on the dual lattice $\Lambda=\operatorname{Hom}(V,\ZZ)$ with dual basis $(\hat\gamma,\hat u,\hat w,\hat\delta)$ \cite{SixSpherePaper}. The package uses the induced affine action on the hyperplane
\[
H_D=\{\lambda\in\Lambda/D\Lambda:\ \text{the $\hat\gamma$-coordinate of $\lambda$ is }1\}\cong(\ZZ/D)^3,
\]
and builds from its orbits covers of the projective line branched at three points.
It presents the cover as a complete failure object: the equation fails at $p$ if and only if a denominator-labelled cover has no section. That equivalence is correct, and it is Theorem~\ref{thm:shell} in cover-theoretic form: a section exists exactly when the tail denominator $R_a/\gcd(R_a,4u+1)$ or $R_a/\gcd(R_a,u+a)$ equals $1$.

For odd levels $D$, which is the package's standing hypothesis and all that the application to the equation uses (there $D$ divides the odd number $R_a$), the package's orbit and genus statements are correct. The action is transitive when $3\nmid D$, and the cover is then connected of genus $(5D^3-12D^2-17D+24)/24$. When $3\mid D$ it has two orbits, of sizes $D^3/9$ and $8D^3/9$, which give components of genera $(5D^3-12D^2-81D+216)/216$ and $(5D^3-12D^2-9D+27)/27$. At even levels these statements do not hold as stated: at $D=2$ the orbits have sizes $2$ and $6$ (found by the referee of the audit notes).
\status{Audited: the proof read; the relations, the orbits and the genera recomputed from cycle counts for $D=1,3,\dots,15$. The equivalence with solvability is a re-encoding of Theorem~\ref{thm:shell}.}

The orbits at every level are determined by one invariant. Write $\lambda=\hat\gamma+b\hat u+c\hat w+d\hat\delta\in H_D$ and $e=\gcd(D,6)$.

\begin{theorem}[the orbits at every level]\label{thm:orbits}
Let the monodromy group $G=\langle T_1,T_2\rangle$ act on $H_D$ dually, through $A_1,A_2$. For every $D\ge1$ its orbits coincide with those of the full stabilizer $G_{\ZZ}$ of the isotropic vector $\gamma$ in the symplectic group of the period lattice, acting in the same way. The orbit of $\lambda$ is
\[
\{\lambda'\in H_D:\ \gcd(b',c',e)=\gcd(b,c,e)\},
\]
of size $D^3\cdot\#\{v\in(\ZZ/e)^2:\ \gcd(v,e)=\gcd(b,c,e)\}/e^2$.
\end{theorem}
\status{New in version~4. Proved in the companion paper on the $S^6$ period system \cite[Theorem~6.1]{SixSpherePaper}, which describes $G$ as the kernel of a character $\Phi$ of order $12$ on $G_{\ZZ}$; the orbit argument is repeated below. Checked as full orbit partitions for every $D\le14$ by the referee script \note{ver52\_referee\_claims.py}, for every $D\le24$ by the referee's own scripts, and for $D=2^k$, $k\le5$, by \note{es09\_even\_levels.py}.}
\begin{proof}
We use from \cite[§§2--3]{SixSpherePaper}: every element of $G_{\ZZ}$ is uniquely $L_Mh(\alpha,\beta,s)$ with $M\in SL_2(\ZZ)$ and $h(\alpha,\beta,s)$ in a Heisenberg group of integer triples; $G$ is the kernel of the homomorphism $\Phi(L_Mh(\alpha,\beta,s))=\operatorname{ab}(M)+s-6(\alpha^2+\alpha\beta+\beta^2)\bmod12$, where $\operatorname{ab}$ is the abelianisation of $SL_2(\ZZ)$ onto $\ZZ/12$ with $\operatorname{ab}\begin{psmallmatrix}1&1\\0&1\end{psmallmatrix}=1$; and $G_{\ZZ}=G\langle t_\gamma\rangle$, where the transvection $t_\gamma=h(0,0,1)$ is central with $\Phi(t_\gamma)=1$.

\emph{The dual action.} The dual action of $h(\alpha,\beta,s)$ is $(b,c,d)\mapsto(b+6\beta,\,c-6\alpha,\,d-s-\alpha b-\beta c)$, and $L_M$ acts on $(b,c)$ by $M^{-\mathsf T}$ and fixes $d$.

\emph{The orbits of $G_{\ZZ}$.} The coordinate $d$ can be moved freely, by $s$, and $(b,c)$ matters only modulo $6(\ZZ/D)^2$. Since $(\ZZ/D)^2/6(\ZZ/D)^2\cong(\ZZ/e)^2$ and $SL_2(\ZZ)$ maps onto $SL_2(\ZZ/e)$, the $G_{\ZZ}$-orbits correspond to the $SL_2(\ZZ/e)$-orbits on $(\ZZ/e)^2$, which are classified by the content $\gcd(v,e)$.

\emph{$G$ has the same orbits.} Since $G$ is normal and $G_{\ZZ}=G\langle t_\gamma\rangle$ with $t_\gamma$ central, it suffices to show $t_\gamma\lambda\in G\lambda$ for every $\lambda$; if this holds at $\lambda$, it holds at $k\lambda$ for every $k\in G_{\ZZ}$, because $t_\gamma k\lambda=kt_\gamma\lambda\in kG\lambda=Gk\lambda$. Every point is $G_{\ZZ}$-equivalent to one with $c=0$, because $SL_2(\ZZ)$ maps onto $SL_2(\ZZ/D)$, which moves every vector of $(\ZZ/D)^2$ to one of the form $(g,0)$. At a point with $(b,c)=(g,0)$, take $M=(T^{\mathsf T})^{-1}$ with $T=\begin{psmallmatrix}1&1\\0&1\end{psmallmatrix}$. Then $M^{-\mathsf T}=T$ fixes $(g,0)$, so $L_M$ fixes $\lambda$; and $M=ST S^{-1}$ with $S=\begin{psmallmatrix}0&-1\\1&0\end{psmallmatrix}$, so $\Phi(L_M)=\operatorname{ab}(T)=1$. Hence $t_\gamma L_M^{-1}$ lies in $G$ and moves $\lambda$ exactly as $t_\gamma$ does.
\end{proof}

\begin{corollary}\label{cor:orbits}
For $D=2^k$, $k\ge1$, there are exactly two orbits, of sizes $D^3/4$ and $3D^3/4$. For odd $D$ there is one orbit if $3\nmid D$, and there are two, of sizes $D^3/9$ and $8D^3/9$, if $3\mid D$. For $6\mid D$ there are four orbits; at $D=6$ their sizes are $6$, $18$, $48$ and $144$, and at $D=12$ they are $48$, $144$, $384$ and $1152$. In general the orbits at $D=2^km$ with $m$ odd are the products of those at $2^k$ and at $m$.
\end{corollary}
\begin{proof}
Count the vectors of $(\ZZ/e)^2$ by content. For $e=2$ the contents $2$ and $1$ occur $1$ and $3$ times; for $e=3$, $1$ and $8$ times; for $e=6$, the contents $6,3,2,1$ occur $1,3,8,24$ times. Multiply by $D^3/e^2$. The product statement is the Chinese remainder theorem, since $e=\gcd(2^k,2)\gcd(m,3)$.
\end{proof}

At the levels $D=2^k$ with $1\le k\le5$ the genera of the two components are $(0,0)$, $(1,3)$, $(17,53)$, $(177,537)$ and $(1569,4721)$ (found by the referee pass of version~1 and recomputed here from cycle counts). By Theorem~\ref{thm:orbits} the components of these covers at every level are indexed by the invariant $\gcd(b,c,\gcd(D,6))$; their genera at all levels are Question~\ref{q:genera}.
